% Figura V9: tabuleiro 4x4 com a solução V = (3,1,4,2)
% Convenção: coluna 1..4 (esq->dir), linha 1..4 (baixo->cima).
\begin{tikzpicture}[scale=1.05, line join=round]
  % fundo das células
  \fill[btFundo] (0,0) rectangle (4,4);
  % grade
  \draw[btNeutro, thick] (0,0) grid (4,4);
  \draw[btNeutro, very thick] (0,0) rectangle (4,4);

  % rótulos de coluna (C1..C4) e linha (L1..L4)
  \foreach \c in {1,2,3,4}{
    \node[font=\scriptsize, text=btNeutro] at (\c-0.5, -0.32) {C\c};
  }
  \foreach \r in {1,2,3,4}{
    \node[font=\scriptsize, text=btNeutro] at (-0.32, \r-0.5) {L\r};
  }

  % rainhas em (coluna/linha): (1,3) (2,1) (3,4) (4,2)
  \foreach \c/\r in {1/3, 2/1, 3/4, 4/2}{
    \begin{scope}[shift={(\c-0.5, \r-0.5)}]
      \fill[btDecisao] (0,-0.06) circle (0.24);
      \fill[btDecisao] (-0.24,0.10) -- (-0.24,0.28) -- (-0.12,0.15)
                       -- (0,0.32) -- (0.12,0.15) -- (0.24,0.28)
                       -- (0.24,0.10) -- cycle;
    \end{scope}
  }
\end{tikzpicture}
